% !TeX root = ../main.tex
\section{Comparing two sounds}
\label{sec:metrics}

The distance $D$ in the estimation formula decides what ``sounds alike'' means.
It is used twice: as the \emph{objective} the search minimises, and as the
\emph{score} a benchmark reports. They need not be the same.

\subsection{Parameter metrics}

If the target was made with the same chain the estimator controls, its true
settings are known, and one can report effect-detection accuracy or F1, the
accuracy of the predicted order, or the mean absolute error (MAE) of knob values
on a 0--1 scale. These need ground truth and say little about the sound: because
different settings can sound alike, a large parameter error can come with a good
sound. Peladeau and Peeters found that training on an audio loss gives better
reconstructions than training on parameters, even when the parameter error is
larger~\cite{peladeau2024blind}. Parameter error should therefore be a
diagnostic, not the main criterion.

\subsection{Aligned audio metrics}

These compare a render $\hat y$ with the target $y$ sample by sample or frame by
frame, so they need the \emph{same performance} on both sides.

\begin{description}[style=nextline,leftmargin=1.2cm]
\item[Error-to-signal ratio (ESR)]
  $\displaystyle \mathrm{ESR} = \frac{\sum_n (y[n]-\hat y[n])^2}{\sum_n y[n]^2}$,
  usually after a pre-emphasis filter that weights high frequencies the way
  hearing does~\cite{wright2020perceptual}. The standard loss for training amp
  models; 0 is perfect. Very sensitive to phase.
\item[Multi-resolution STFT distance (MR-STFT)] The difference of magnitude
  spectrograms at several window sizes (e.g.\ 512, 1024, 2048), as a sum of a
  spectral-convergence term and the L1 distance of log
  magnitudes~\cite{yamamoto2020parallel,steinmetz2020auraloss}. Ignores phase;
  the several resolutions catch both fast transients and fine pitch detail.
\item[SI-SDR] Scale-invariant signal-to-distortion ratio: project $\hat y$ on
  $y$ to get the best-scaled target $\alpha y$, then
  $10\log_{10}\bigl(\|\alpha y\|^2/\|\hat y - \alpha y\|^2\bigr)$, in dB, higher
  is better~\cite{leroux2019sdr}.
\item[Log-band distance (the app's)] Mean absolute difference, in dB, between
  log energies in a set of frequency bands, over frames; defined exactly in
  Section~\ref{sec:app-distance}. Close in spirit to the log-magnitude part of
  MR-STFT.
\end{description}

\subsection{Alignment-free metrics}

These compare clips without requiring the same notes:
\begin{itemize}
\item \textbf{Effect embeddings}: networks trained to describe the processing
  independently of the content, such as AFx-Rep~\cite{steinmetz2024stito} and
  Fx-Encoder++~\cite{yeh2025fxencoder}, or general audio embeddings such as
  CLAP~\cite{wu2023laionclap}. Two clips are compared by the cosine distance of
  their embeddings.
\item \textbf{Distribution distances}: the Fréchet Audio Distance
  (FAD)~\cite{kilgour2019fad} fits a Gaussian to the embeddings of many clips of
  each set and measures how far apart the two Gaussians are. It is biased by
  sample size and depends on the embedding and the reference set; FAD$_\infty$
  extrapolates the bias away~\cite{gui2024adapting}, and the kernel variant KAD
  converges with fewer samples~\cite{chung2025kad}. A single song gives few clips,
  so per-song evaluation suits pairwise embedding distances better than FAD.
\end{itemize}

\subsection{Listening tests}

The reference is human judgement. \term{MUSHRA} (ITU-R BS.1534-3)~\cite{itu2015bs1534}
asks listeners to rate several versions against a reference on a 0--100 scale,
with a hidden copy of the reference and a degraded anchor as checks;
webMUSHRA~\cite{schoeffler2018webmushra} runs it in a browser. \term{Mean opinion
scores} (1--5) can be split into tone similarity and audio
quality~\cite{yeh2026egvae}. As a ceiling, listeners could barely tell Kemper
profiles from the original amplifiers ($d' = 0.34$, $N = 177$)~\cite{duvel2020confusingly}.

\subsection{Which metrics agree with listeners}

Two studies measured this for guitar devices (Table~\ref{tab:metrics}). In the
ToneTwist evaluation of 40 nonlinear devices (27 screened listeners), MR-STFT
had the strongest correlation with MUSHRA ratings among signal metrics and ESR
one of the weakest; among FAD embeddings, VGGish and AFx-Rep did best overall
and CLAP best on distortion; no metric was reliable for overdrive or compression
\reported~\cite{comunita2025differentiable}. Wright, Dal Santo and Schlecht
collected ratings for 590 emulation--reference pairs and found MFCC and deep
embedding distances predictive, MAE and ESR poor \reported~\cite{wright2026predicting}.
In unpaired amp modelling, ESR and listener ratings disagreed on a heavily
distorted tone~\cite{wright2023adversarial}. No metric has been validated for
delay, reverb or modulation on guitar.

\begin{table}[h]
\centering
\caption{Spearman correlation of objective metrics with MUSHRA ratings on 40
nonlinear guitar and bass devices~\cite{comunita2025differentiable}. Negative:
lower distance goes with higher rating, so more negative is better.}
\label{tab:metrics}
\small
\begin{tabular}{@{}lrlr@{}}
\toprule
Signal metric & $\rho$ & FAD embedding & $\rho$ \\
\midrule
MR-STFT & $-0.427$ & VGGish & $-0.430$ \\
MAPE & $-0.339$ & AFx-Rep & $-0.417$ \\
L1 & $-0.332$ & CLAP & $-0.406$ \\
MSE & $-0.314$ & PANN & $-0.303$ \\
ESR & $-0.311$ & & \\
\midrule
\multicolumn{4}{@{}l@{}}{High-distortion devices only: MR-STFT $-0.536$, FAD with CLAP $-0.613$.} \\
\bottomrule
\end{tabular}
\end{table}

\begin{caution}[: correlations are modest]
The best correlation in Table~\ref{tab:metrics} is about $-0.43$: objective
metrics explain only part of what listeners hear. A lower distance is evidence
of a better match, not proof. Listening tests are on the app's roadmap
(Section~\ref{sec:next}); the benchmark already writes WAV files for listening
(\code{GVA_TONE_OUT}).
\end{caution}

\subsection{Comparability of published results}

Most results use clean, isolated, plugin-rendered guitar, which is easier than
the commercial case. Many test sets share dry recordings or plugins with the
training set. And because settings can be non-unique, two estimates with very
different parameters may sound alike, which parameter metrics penalise and
listeners do not.
